\section*{Introduction}

Natural environments are richly structured by statistical regularities. Statistical learning (SL) refers to the set of processes by which learners detect, internalize, and exploit such regularities in sensory input. Frost et al.\ define SL as “perceiving and learning any forms of patterning in the environment that are either spatial or temporal in nature,” and using this information to construct mental representations (p.~1130) \cite{frostStatisticalLearningResearch2019}. Interest in SL has grown substantially within cognitive neuroscience because sensitivity to structured regularities supports many core cognitive functions, including contextual cueing, scene perception, visual word recognition, and face perception \cite{cleeremansImplicitLearningNews1998, frostStatisticalLearningResearch2019, daltrozzoNeurocognitiveMechanismsStatisticalsequential2014, shermanPrevalenceImportanceStatistical2020}. Reading acquisition, for example, depends on learning probabilistic regularities linking orthography, phonology, morphology, and syntax, as well as on sensitivity to the co-occurrence structure of letters and words \cite{frostStatisticalLearningResearch2019}.

Although SL is often assumed to operate implicitly, much of the empirical literature has relied on behavioral measures that require explicit judgments about what has been learned. Such measures typically reflect the end point of multiple processing stages, including perceptual encoding, decision-making, response strategies, and task demands. Consequently, learning-related representations may be present at a neural level even when behavioral performance is weak, variable, or inconsistent. This possibility motivates the use of neural measures to complement behavioral approaches; such measures can reveal sensitivity to statistical structure even when behavioral performance is weak or unreliable.

Event-related potentials (ERPs) are particularly well suited to this goal. ERPs provide millisecond-level temporal resolution and are time-locked to stimulus onset, allowing learning-related effects to be traced across distinct stages of information processing. In addition, ERP components can be differentiated by their latency and scalp distribution, enabling inferences about the functional locus of statistical learning effects within the perceptual–cognitive hierarchy \cite{curranBrainPotentialsRecollection2000}. For these reasons, ERP methods have been widely used to investigate the neural dynamics of statistical learning across sensory modalities. Critically, the temporal precision allows ERPs to reveal when statistical structure influences processing, even in the absence of explicit behavioral evidence of learning.

In the present study, we use ERPs to examine neural sensitivity to statistical structure in a visual statistical learning paradigm. We focus on two components that have been repeatedly implicated in prior SL research: the N100, associated with early perceptual and attentional processes, and the N400, associated with later integration and representational updating. Although both components have been linked to statistical learning in auditory paradigms, their respective roles in visual statistical learning—particularly with respect to early perceptual segmentation and the expression of learned structure—remain less well understood.

\subsection*{The N100 and N400 Components}

The N100 is an early negative-going ERP component that typically peaks between approximately 80 and 120~ms following stimulus onset. In auditory paradigms, it is usually maximal over fronto-central scalp sites, whereas in visual paradigms it is most prominent over posterior regions. The N100 is generally interpreted as reflecting early perceptual encoding of stimulus features and is strongly modulated by attention, salience, and expectancy \cite{luckIntroductionEventRelatedPotential2014, hillyardEventrelatedBrainPotentials1998}. Attended or behaviorally relevant stimuli reliably elicit larger N100 amplitudes than unattended stimuli, indicating that this component indexes early, attention-sensitive stages of sensory processing \cite{hillyardEventrelatedBrainPotentials1998, correaAttentionalMechanismTemporal2006}. As such, the N100 provides a useful marker of how learned regularities can influence perceptual processing at an early stage.

The N400 is a later negative-going component that typically peaks between 300 and 500~ms post-stimulus and is maximal over centro-parietal scalp sites. Originally identified in studies of language comprehension, the N400 was linked to semantic processing and the integration of words into sentence context \cite{kutasReadingSenselessSentences1980}. Subsequent research has demonstrated that N400 amplitude is sensitive more broadly to predictability, associative structure, and the ease with which incoming input can be integrated into an existing representational framework \cite{kutasThirtyYearsCounting2011}. In this broader view, the N400 reflects a general process of integration or updating, indexing the degree to which a stimulus coheres with learned structure rather than semantic content per se. This broader interpretation makes the N400 a particularly useful index of statistical learning, where structure must be inferred from distributional regularities rather than semantic content.


\subsection*{ERP Evidence for Statistical Learning}

A substantial body of ERP research has identified the N400 as a robust neural marker of statistical learning, particularly in auditory paradigms. In these studies, participants are typically exposed to continuous streams of syllables or tones in which statistically defined units emerge through high transitional probabilities between elements. Across a wide range of task designs, the N400 reliably differentiates structured from unstructured input, indicating sensitivity to distributional regularities in the absence of explicit segmentation cues.

Early work by Sanders et al.\ \cite{sandersSegmentingNonsenseEventrelated2002} demonstrated that N400 amplitude increased following exposure to statistically coherent syllable streams, even among participants who showed little or no behavioral evidence of learning. Subsequent studies have shown that N400 responses evolve systematically over the course of exposure, often exhibiting an inverted-U trajectory in which amplitudes increase as structure is detected and integrated, then decrease as representations become more stable and familiar \cite{ablaOnlineAssessmentStatistical2008, cunilleraTimeCourseFunctional2009}. This pattern suggests that the N400 reflects the formation and consolidation of internal representations based on statistical structure, rather than behavioral success or explicit recognition per se.

In contrast, evidence for N100 modulation in statistical learning has been more variable and appears to depend on the availability of cues that facilitate early segmentation or attentional orienting. In auditory SL paradigms, N100 enhancement has been reported primarily when strong boundary cues are present, such as prosodic stress, highly predictive transitional probabilities, or explicit instructions directing attention to structure \cite{cunilleraEffectsStressStatistical2006, sandersERPStudyContinuous2003, soaresNotAllWords2020, soaresLearningWordsListening2022}. Under these conditions, the N100 has been interpreted as reflecting early perceptual segmentation or attentional allocation to statistically defined units. When such cues are absent, however, N100 effects are often weak or absent, even in the presence of robust N400 modulation, underscoring a dissociation between early perceptual segmentation and later integrative processes in statistical learning. Together, these findings suggest that early perceptual markers of statistical learning are more contingent on task structure and cue availability than later integration-related processes.

\subsection*{Visual Statistical Learning and Open Questions}

ERP investigations of statistical learning in the visual domain have been comparatively rare, and their findings have been less consistent than those reported in auditory paradigms. Many visual SL studies have relied on target-detection or contingency-learning designs that emphasize associative relationships between stimuli rather than the segmentation of continuous sequences. As a result, these paradigms are not well suited to probing early perceptual components or to directly comparing the neural dynamics of learning across sensory modalities, and it remains unclear whether the absence of early visual effects reflects a true modality difference or limitations of existing experimental designs.

In that study, participants viewed continuous streams of geometric shapes organized into repeating triplets and were later tested on their ability to discriminate familiar from unfamiliar sequences. ERP responses were recorded during exposure to the stream. The authors observed N400 modulation at triplet boundaries among high-performing learners, suggesting that visual statistical learning engages later-stage integration mechanisms similar to those observed in auditory SL. However, no modulation of the N100 component was observed, even among participants who demonstrated behavioral learning. Moreover, low-performing learners did not show reliable N400 effects, in contrast to many auditory SL studies in which N400 modulation has been observed despite weak behavioral performance.

These findings raise two unresolved questions. First, does visual statistical learning reliably engage early perceptual or attentional mechanisms indexed by the N100, or is such modulation contingent on task design, learning stage, or the presence of explicit segmentation cues? Second, can neural evidence of learning---particularly N400 modulation---be observed in visual SL even when behavioral discrimination is weak or inconsistent, as has been reported in auditory paradigms? Addressing these questions is critical for understanding whether the neural dynamics of statistical learning are truly domain general or are shaped by modality-specific constraints. Answering these questions requires a design that can isolate neural sensitivity to learned structure independently of online segmentation demands and overt behavioral success.

\subsection*{The Present Study}

Critically, unlike prior visual SL studies, the present study addresses these open questions by examining ERP responses during a post-familiarization test phase in a visual statistical learning paradigm. Participants were exposed to continuous sequences of shape triplets without explicit segmentation cues and were subsequently tested on their ability to discriminate familiar from unfamiliar triplets. ERPs were recorded during the test phase, allowing us to assess neural sensitivity to learned structure after acquisition, rather than during online learning as in prior visual SL studies (e.g., Abla \& Okanoya, 2009). This design allows learned structure to shape neural responses prospectively, rather than requiring it to emerge online during exposure.

Using a larger sample ($N = 67$) than previous ERP studies of visual statistical learning, our primary goal was to determine whether neural responses reliably differentiate familiar from unfamiliar sequences at the group level. In addition, we examined whether neural familiarity effects were modulated by trial-level behavioral accuracy during the familiarity judgment task. This approach allowed us to test whether neural indices of learning covary with overt performance or instead reflect sensitivity to statistical structure that is expressed independently of explicit decision outcomes.

Accordingly, we addressed two central questions: (1) Do the N100 and N400 components reliably differentiate familiar from unfamiliar visual sequences following statistical learning? and (2) Are these neural familiarity effects modulated by trial-level accuracy on the behavioral test?

Based on prior auditory and visual SL research, we predicted that the N400 would show robust familiarity-related modulation, reflecting the integration of learned structure. Given the mixed and inconsistent evidence for early perceptual effects in visual SL, predictions regarding N100 modulation were treated as exploratory. By combining a well-powered design with linear mixed-effects modeling, the present study aims to clarify when and how early and late ERP components index statistical learning in the visual domain and to dissociate neural sensitivity to structure from explicit behavioral performance.


\section*{Materials and methods}

\subsection*{Participants}
We tested 76 participants recruited from the UMass–Five Colleges consortium community in Amherst, MA between September 2017 and May 2019. All procedures were approved by the Institutional Review Board at Hampshire College. The IRB does not issue protocol numbers; a formal approval letter is included as supporting documentation. Participants were healthy adult volunteers who provided written informed consent prior to participation, in accordance with the Declaration of Helsinki. Identifying information was not collected, and all data were analyzed in de-identified form.

Participants were eligible if they were between the ages of 18 and 26, right-handed, native English speakers with no other languages spoken in the home before age 3, and reported normal or corrected-to-normal vision with no history of cognitive or reading-related difficulties. Three participants were excluded for exceeding the protocol’s age range, and six were excluded due to data loss, resulting in a final sample of 67 participants (age range: 18–25 years, $M = 19.6$). Of these, 49 identified as female, 17 as male, and one did not report their sex. Participants were compensated for their time.

To assess whether the study was adequately powered to detect group-level effects of a magnitude comparable to those reported in prior ERP studies of statistical learning, we conducted a post hoc power analysis using effect sizes reported in the literature. Published $\eta^2$ values for key familiarity-related effects typically range from .25 to .44. Converting these values to Cohen’s $f$ and applying them to our final sample size ($n = 67$), we estimated high power ($\geq .95$) to detect effects of similar magnitude at the group level. As with many ERP studies, the present design was not optimized to detect higher-order interactions or individual-difference effects, which are interpreted descriptively.

\subsection*{Stimuli}
The stimuli and procedure were adapted from Siegelman and Frost (2015) \cite{siegelmanStatisticalLearningIndividual2015}, which in turn were based on Glicksohn and Cohen (2013) \cite{glicksohnRoleCrossmodalAssociations2013}. During an initial familiarization phase, participants viewed a continuous stream of unfamiliar visual shapes presented at the center of the screen against a white background. The stimulus set consisted of 24 unique shapes, organized into eight visual triplets. Each triplet comprised three distinct shapes that always appeared consecutively, with triplet assignment randomized separately for each participant.

The familiarization stream lasted approximately 10 minutes and consisted of repeated presentations of the eight triplets in random order, with the constraint that the same triplet could not appear twice in succession. Each triplet was presented 24 times during familiarization. Individual shapes were displayed for 800 ms, separated by a 200 ms inter-stimulus interval. Participants were instructed simply to attend to the shapes and were not informed about the presence of triplet structure.

Following familiarization, participants completed a two-alternative forced-choice test designed to assess learning of the triplet structure. On each trial, participants viewed two sequentially presented triplets: one familiar triplet from the training phase and one foil triplet composed of three shapes that had never appeared together during familiarization. Participants indicated which triplet seemed more familiar by pressing a button on the keyboard. For familiar triplets, transitional probabilities between successive shapes were always 1.0, whereas for foil triplets transitional probabilities were 0. The test comprised 32 trials, yielding performance scores ranging from 0 to 32. Participants completed the test after finishing a language battery and an unrelated lexical decision task.

\subsection*{EEG Data Acquisition and Processing}
EEG activity was recorded at 500 Hz from 31 Ag/AgCl active electrodes (actiCAP; BrainProducts GmbH, Gilching, Germany) placed according to the extended 10--20 system and mounted in an elastic cap. Data were recorded using an initial reference-free montage and were re-referenced offline to the algebraic average of the left and right mastoid electrodes. The ground electrode was placed at Fpz. Electrode impedances were kept below 25~$k\Omega$, which is within the acceptable range for active electrodes. Signals were amplified using an actiCHamp Plus EEG system and recorded with BrainVision Recorder software (BrainProducts GmbH). Ocular activity (blinks and eye movements) was monitored using frontal electrodes (FP1 and FP2) and corrected using ICA.

EEG data were preprocessed using EEGLAB \cite{delormeEEGLABOpenSource2004} and ERPLAB \cite{lopez-calderonERPLABOpensourceToolbox2014}. The continuous data were bandpass filtered from 0.1 to 30 Hz using a zero-phase Butterworth filter, downsampled to 200 Hz, and re-referenced to the algebraic average of the left and right mastoid electrodes. Independent component analysis (ICA) was performed to identify and remove components reflecting ocular and muscular artifacts. Components associated with eye blinks and saccades were identified based on their characteristic topography and time course and were removed prior to further processing. Bad electrodes were identified and interpolated using spherical interpolation.

The continuous data were segmented into epochs extending from –200 ms to 1000 ms relative to the onset of each shape. Baseline correction was applied using the 100 ms pre-stimulus interval (–100 to 0 ms relative to stimulus onset), a window commonly used in ERP research to normalize voltage fluctuations and isolate stimulus-evoked activity. Epochs containing voltage fluctuations exceeding $\pm100~\mu$V were excluded from further analysis. After artifact rejection, 96.5\% of trials were retained.

Epoched data were averaged to create event-related potentials (ERPs). Mean amplitudes were extracted from two post-stimulus time windows corresponding to the N100 (75--175 ms) and N400 (350--500 ms). These intervals were selected based on prior ERP research examining early perceptual and later integration-related components in visual statistical learning and visual object processing paradigms.


\subsection*{Sensitivity Classification}
Participants’ performance on the two-alternative forced-choice familiarity test was used to index between-participant variability in learning strength. Behavioral performance was defined as the number of correct responses on the 2AFC test (range: 0--32), and proportion correct ($p_c$) was computed as correct/32. For descriptive purposes, we also expressed 2AFC performance in signal-detection units using the standard equal-variance conversion for two-alternative forced choice, $d' = \sqrt{2}\, z(p_c)$, where $p_c$ is the proportion correct. Because this sensitivity metric is a monotonic transformation of accuracy, it was not treated as an independent outcome variable but as a descriptive summary of discrimination performance. Participants with $d' > 0$ (equivalently, $p_c > .50$) were classified as sensitive, whereas those with $d' \leq 0$ were classified as insensitive. Based on this criterion, 17 participants (25.4\%) were categorized as sensitive and 50 participants (74.6\%) as insensitive. Mean accuracy was higher for sensitive participants (79.3\%) than for insensitive participants (51.7\%), reflecting the shared variance between these measures.

\subsection*{Statistical Analysis}
Statistical analyses were conducted using linear mixed-effects modeling in R version 4.4.2 \cite{rcoreteamLanguageEnvironmentStatistical2024} within RStudio \cite{rstudioteamRStudioIntegratedDevelopment2024}, using the \textit{afex} package \cite{singmannAfexAnalysisFactorial2024}. A mixed-effects approach was selected over traditional repeated-measures ANOVA because it accommodates unbalanced data, accounts for subject-level variability, and does not require sphericity assumptions, making it well suited to ERP data with multiple repeated measurements per participant \cite{tremblayModelingNonlinearRelationships2015, barrRandomEffectsStructure2013, mageziLinearMixedeffectsModels2015}.

\paragraph{Primary (group-level) analyses}
Primary analyses focused on group-level effects of \textit{Familiarity} (Familiar, Unfamiliar; within-subjects), independent of individual differences in behavioral performance. These models were used to establish whether neural responses differentiated familiar from unfamiliar triplets at the population level. Sensitivity was not included in the primary models because it is derived from behavioral accuracy and therefore shares variance with trial-level correctness.

ERP amplitude was modeled without additional fixed effects for scalp location (e.g., Anteriority or Laterality), as these spatial factors were not central to the hypotheses and including them would substantially increase the number of statistical comparisons, inflating the Type I error rate \cite{luckHowGetStatistically2017}. To account for spatially correlated observations, electrode site was instead modeled as a random intercept nested within participant. The primary model was specified as:

\begin{eqnarray}
	\texttt{ERP amplitude} &\sim& 
	\texttt{Familiarity} \nonumber \\
	&& +\ \texttt{(1 + Familiarity | SubjectID)} \nonumber \\
	&& +\ \texttt{(1 | SubjectID:Electrode Site)}
\end{eqnarray}

\paragraph{Secondary (individual-differences) analyses}
To characterize how group-level familiarity effects were distributed across individuals, secondary analyses included \textit{Sensitivity} (Sensitive, Insensitive; between-subjects) and its interaction with Familiarity. Because Sensitivity is derived from behavioral accuracy, these models were interpreted descriptively and were not used as the basis for confirmatory inference. Instead, they were used to assess whether the magnitude or direction of familiarity-related neural effects varied systematically as a function of individual differences in learning strength.

In these models, random intercepts for each participant captured baseline variability in ERP amplitude, while random slopes for Familiarity allowed for individual differences in the size and direction of familiarity effects. Electrode sites were nested within participants to retain spatial variability while avoiding the multiple comparisons and inflated familywise error rate associated with modeling scalp region as a fixed effect \cite{luckHowGetStatistically2017}. This approach reduces the number of hypothesis tests, avoids information loss associated with averaging across electrodes, and more accurately reflects the nested structure of ERP data than a cross-classified model treating subject and electrode as independent random effects.

\paragraph{Model estimation and inference}
Models were estimated using restricted maximum likelihood (REML), and the Kenward--Roger approximation was used to calculate denominator degrees of freedom and improve control of Type I error \cite{kenwardSmallSampleInference1997}. Fixed effects were tested using Type III Wald $F$ tests. Estimated marginal means (EMMs), standard errors (SE), and 95\% confidence intervals (CI) were obtained using the \textit{emmeans} package \cite{lenthEmmeansEstimatedMarginal2024}, and pairwise comparisons were Bonferroni-corrected to account for multiple testing. The full random-effects structure was retained following recommendations for confirmatory hypothesis testing in ERP designs \cite{barrRandomEffectsStructure2013}.



\section*{Results}

To address our research questions, we analyzed ERP mean amplitudes in two time windows: the N100 (75--175 ms), indexing early perceptual processing, and the N400 (350--500 ms), often linked to later integration-related processing. Our primary aim was to test whether ERP responses reliably differentiated statistically familiar from unfamiliar shape sequences. A secondary aim was to assess whether ERP effects covaried with trial-level response accuracy on the familiarity test. In all models, electrode site was modeled as a random intercept nested within participant, and by-participant random slopes were included for within-subject predictors (see Statistical Analysis).

\subsection*{Primary analyses: Familiarity and Accuracy}

For each component, we fit a linear mixed-effects model with fixed effects of Familiarity (familiar vs.\ unfamiliar), Accuracy (correct vs.\ incorrect), and their interaction, with random slopes for Familiarity and Accuracy by participant and random intercepts for electrode site nested within participant.

\paragraph{N100 (75--175 ms).}
The model revealed a significant main effect of Familiarity, $F(1,66) = 29.62$, $p < .001$, $\eta_p^2 = .31$, with more negative N100 amplitudes for familiar than unfamiliar sequences. Neither the main effect of Accuracy, $F(1,66) = 1.05$, $p = .310$, nor the Familiarity $\times$ Accuracy interaction, $F(1,1674) = 0.78$, $p = .377$, was significant. Estimated marginal means (averaged over Accuracy) indicated a more negative N100 for familiar sequences (M = $-3.18~\mu$V, SE = 0.44, 95\% CI [$-4.07$, $-2.30$]) than for unfamiliar sequences (M = $-0.78~\mu$V, SE = 0.23, 95\% CI [$-1.23$, $-0.33$]), yielding a mean difference of $-2.40~\mu$V, $p < .0001$, $d_z = -0.55$. This pattern indicates that learned statistical structure influenced early perceptual processing, independent of whether participants responded correctly on a given trial.

\paragraph{N400 (350--500 ms).}
The corresponding model for the N400 also showed a robust main effect of Familiarity, $F(1,66) = 36.95$, $p < .001$, $\eta_p^2 = .36$, again reflecting more negative amplitudes for familiar than unfamiliar sequences. The main effect of Accuracy was not significant, $F(1,66) = 0.82$, $p = .369$. The Familiarity $\times$ Accuracy interaction showed a non-significant trend, $F(1,1674) = 2.98$, $p = .085$. Estimated marginal means (averaged over Accuracy) indicated a more negative N400 for familiar sequences (M = $-2.00~\mu$V, SE = 0.55, 95\% CI [$-3.09$, $-0.91$]) than for unfamiliar sequences (M = $1.77~\mu$V, SE = 0.50, 95\% CI [$0.77$, $2.77$]), corresponding to a mean difference of $-3.77~\mu$V, $p < .0001$, $d_z = -0.55$. Thus, later integration-related neural responses were robustly shaped by familiarity with statistical structure, even when explicit discrimination was inconsistent.

Together, these primary analyses demonstrate that both the N100 and N400 reliably differentiated familiar from unfamiliar shape sequences. Critically, this neural differentiation was present independent of trial-level response accuracy on the behavioral familiarity test, indicating that ERP indices of learning were robust at the group level even when behavioral discrimination performance varied across trials.

Finally, to assess whether individual differences in behavioral sensitivity predicted the magnitude of neural familiarity effects, we correlated participant-specific random slopes for Familiarity with behavioral sensitivity ($d’$). This analysis revealed no reliable association between neural and behavioral measures ($r = .15$, $p = .23$, 95\% CI [$-.09$, $.38$]), indicating that variability in ERP familiarity effects was not systematically related to explicit discrimination performance.



\section*{Discussion}
This study investigated how neural responses to statistical structure in visual input reflect prior exposure to shape sequences and the extent to which these responses align with behavioral performance. Using event-related potentials (ERPs), we focused on two components—the N100 and N400—that have previously been associated with early perceptual and attentional processes and later integration-related processing, respectively. Our investigation was guided by two central questions: First, do the N100 and N400 components reliably distinguish between shape sequences that were previously encountered during familiarization and those that were novel, indicating neural sensitivity to learned structure? Second, to what extent do these neural responses covary with individual differences in behavioral performance on a post-exposure familiarity judgment task? Our findings provide clear answers to both questions and clarify the relationship between neural indices of statistical learning and overt behavioral discrimination. In short, the brain distinguished familiar from unfamiliar visual sequences even when participants often could not.

Both the N100 and N400 components showed robust modulation by Familiarity, with significantly more negative amplitudes for familiar compared to unfamiliar sequences. These findings demonstrate that statistical learning influences neural responses at both early perceptual (N100) and later integration-related (N400) stages of processing. The presence of strong familiarity effects in both components, with substantial effect sizes, aligns with prior work in auditory statistical learning and extends these findings to the visual modality in a well-powered sample.

A central finding of the present study is that robust neural differentiation between familiar and unfamiliar sequences was observed even when behavioral discrimination performance was weak or inconsistent. Across participants, familiarity-related ERP effects emerged independently of trial-level response accuracy and did not scale reliably with individual differences in behavioral sensitivity. This dissociation suggests that neural indices of statistical learning can reflect the extraction of regularities from the input even when this knowledge is not consistently accessible to conscious report or expressed in overt task performance. In this respect, the present findings diverge from several auditory statistical learning studies in which stronger behavioral learners often exhibit larger neural responses to structured sequences \cite{sandersSegmentingNonsenseEventrelated2002, ablaOnlineAssessmentStatistical2008}.

The lack of a systematic relationship between neural familiarity effects and behavioral sensitivity highlights an important distinction between learning-related neural representations and the processes that support explicit discrimination judgments. While behavioral performance reflects decision-making under task demands, ERP measures capture online neural responses to structured input that may operate independently of conscious awareness or response strategies. The present results therefore support the view that ERPs can index implicit learning processes that are not fully captured by overt behavioral measures, particularly in tasks that require explicit familiarity judgments.

The remainder of the discussion is organized around two central questions raised by these findings. First, we consider how neural sensitivity to statistical regularities can emerge in the absence of reliable behavioral discrimination, and what this dissociation reveals about implicit learning mechanisms. Second, we examine the conditions under which early N100 effects arise in statistical learning paradigms, with particular attention to why the present study revealed robust N100 modulation in a visual task where previous studies have not. Together, these issues inform a broader account of how statistical structure is encoded, expressed, and accessed across multiple levels of processing.

\subsection*{Neural Evidence of Learning in the Absence of Reliable Behavioral Discrimination}

The robust N400 familiarity effects observed in the present study challenge the assumption that neural indices of statistical learning must closely track overt behavioral performance. Across participants, reliable differentiation between familiar and unfamiliar sequences emerged at the neural level even when behavioral discrimination was weak or inconsistent. This dissociation raises an important question: how can clear neural sensitivity to statistical structure arise in the absence of reliable explicit judgments of familiarity?

One possibility is that neural sensitivity to structure reflects learning processes that precede, or operate independently from, the mechanisms required for explicit behavioral report. In this view, the extraction and representation of statistical regularities can be well established at a neural level even when those representations are not consistently accessed or used during decision-making. Prior longitudinal studies of statistical learning provide support for this account. In a foundational auditory statistical learning study, Abla et al.\ \cite{ablaOnlineAssessmentStatistical2008} recorded ERPs during the exposure phase and showed that N400 responses to structured sequences evolved over time, with changes in amplitude reflecting the gradual stabilization of learned representations rather than moment-to-moment behavioral performance.

These findings are consistent with an inverted-U trajectory in the N400 component during learning. Early in learning, as structure is first detected and representations are unstable, integration demands are high, resulting in larger N400 responses. With continued exposure and consolidation, representations become more stable and easier to integrate, leading to reduced N400 amplitudes. Importantly, even when the N400 decreases, it often remains elevated relative to unstructured input, indicating sustained sensitivity to learned regularities \cite{cunilleraTimeCourseFunctional2009}.

Within this framework, the N400 reflects the effort required to integrate incoming input into an internal model of learned structure. When no structure is present, as in random sequences, there is little basis for prediction or integration, and the N400 is minimal. As statistical regularities emerge but remain fragile, integration becomes effortful, yielding large N400 responses. Once structure is well established, integration is more efficient and the N400 correspondingly attenuates. Thus, N400 amplitude indexes representational instability rather than learning success per se.

This interpretation parallels classic findings from lexical-semantic ERP research. In lexical decision tasks, pronounceable nonwords typically elicit larger N400 responses than real words, while consonant strings evoke little or no N400 \cite{kutasThirtyYearsCounting2011}. Nonwords activate lexical processing mechanisms but lack stable representations, resulting in heightened integration demands. Real words, by contrast, are highly familiar and richly represented, allowing for efficient integration. Consonant strings fail to engage lexical processing altogether. A similar gradient appears to operate in statistical learning paradigms, even in the absence of semantic content.

In the present study, ERPs were recorded during a post-familiarization test phase rather than during exposure. By this point, participants may have acquired statistical regularities to varying degrees, but the accessibility of those representations for explicit familiarity judgments may differ. Robust N400 familiarity effects under these conditions suggest that internal representations of structure were sufficiently developed to support prediction and integration, even when behavioral responses did not reliably reflect that knowledge.

The absence of a systematic relationship between N400 familiarity effects and behavioral sensitivity further supports this interpretation. Participant-specific neural familiarity effects did not scale with behavioral discrimination performance, indicating that neural learning signals were not simply weaker or delayed in participants with lower accuracy. Instead, the results point to a dissociation between the presence of learned structure at a neural level and the ability to deploy that knowledge in an explicit decision-making task.

Finally, the use of linear mixed-effects models may have increased sensitivity to these learning-related neural signals by accounting for both subject- and electrode-level variability. Such models are well suited to detecting effects that are distributed unevenly across trials or participants, and may therefore reveal neural learning effects that are missed by approaches relying on extensive averaging.

Taken together, these findings suggest that N400 familiarity effects in visual statistical learning reflect implicit sensitivity to structure that can emerge independently of explicit behavioral discrimination. Rather than indicating anomalous or failed learning, neural familiarity effects in the absence of reliable behavioral performance highlight the capacity of ERPs to index latent learning processes that are not fully captured by overt responses. These results indicate that neural representations of statistical structure can be robust even when access to those representations for explicit judgment is unreliable.


\subsection*{When Do N100 Effects Arise in Statistical Learning?}

Prior ERP studies of statistical learning (SL) have reported mixed evidence regarding N100 modulation. In some auditory paradigms, enhanced N100 responses have been observed at word or triplet onsets, particularly when external segmentation cues such as prosodic stress \cite{cunilleraEffectsStressStatistical2006, sandersERPStudyContinuous2003} or very strong transitional probability (TP) manipulations \cite{soaresNotAllWords2020, soaresLearningWordsListening2022} are present. These cues are thought to guide attention to likely boundaries, thereby enhancing early perceptual processing. For example, Cunillera et al.\ reported larger N1 and N400 responses to stressed syllables, while Soares and colleagues observed stronger N100 and N400 responses to high-TP items. In contrast, when such cues are absent and structure is defined only by moderately strong TPs, N100 effects are often weak or absent. Consistent with this pattern, Abla and Okanoya \cite{ablaVisualStatisticalLearning2009} reported clear N400 effects but no N100 modulation in a visual SL task without explicit boundary markers.

One possible explanation for the presence of N100 effects in the current study concerns the timing of ERP measurement. Here, ERPs were recorded during a post-familiarization test phase, after participants had substantial opportunity to acquire the statistical structure of the shape sequences. In contrast, several prior studies, including that of Abla and Okanoya \cite{ablaVisualStatisticalLearning2009}, recorded ERPs during the exposure phase, when structure-driven attentional modulation may not yet have emerged. If the N100 reflects attention guided by learned structure, such effects may only become detectable once learning has occurred and when that knowledge can be prospectively deployed. This distinction between online learning and offline recognition may therefore contribute to divergent N100 findings across studies.

A second explanation relates to differences in stimulus presentation and segmentation. In the present study, ERPs were measured during a test phase in which triplets were presented as isolated units with clearly defined boundaries. By contrast, in exposure-phase recordings, triplets typically appear within a continuous stream of unsegmented input. If the N100 reflects a perceptual response to the onset of a learned unit, it may be more readily elicited when that onset is unambiguous, as in isolated trials, than in continuous streams where boundaries must be inferred. On this view, N100 modulation in the present study reflects the interaction between learned structure and clear segmentation at test, rather than reliance on exogenous boundary cues.

Importantly, robust N100 modulation was observed across participants regardless of their level of explicit behavioral discrimination. This pattern suggests that the N100 in the current task does not primarily reflect cue-driven or response-contingent processing, but rather endogenously guided attention shaped by acquired statistical structure. Once learned, statistical regularities may themselves function as attentional cues, supporting early perceptual grouping and enhancing N100 responses. Prior studies that failed to observe N100 effects may therefore have captured neural activity too early in the learning trajectory or under conditions in which learned structure could not yet be effectively deployed to guide perception. Early perceptual markers of statistical learning may therefore reflect the deployment of acquired structure, rather than the initial discovery of that structure.




\section*{Conclusion}

This study provides clear evidence that visual statistical learning modulates neural responses at both early (N100) and later (N400) stages of processing. Across participants, familiar shape sequences elicited robust and reliable ERP differentiation relative to unfamiliar sequences, demonstrating that the brain can track statistical structure in visual input even when overt behavioral discrimination is weak or inconsistent. These effects were observed in a well-powered sample using consistent analytical procedures and convergent measures across components, underscoring the reliability of the findings.

Importantly, neural familiarity effects were largely independent of behavioral performance. Familiarity-related ERP responses did not depend on trial-level accuracy and did not scale systematically with individual differences in behavioral sensitivity. This dissociation suggests that the acquisition and neural representation of statistical regularities can proceed independently of the ability to access or deploy that knowledge during explicit decision-making.

Together, these findings reinforce the value of ERPs as a window into latent learning processes that may not be fully captured by overt behavior. By revealing neural sensitivity to structure in the absence of reliable behavioral expression, the present results highlight the importance of integrating neural and behavioral measures when characterizing learning—particularly in domains where knowledge may be implicit, fragile, or not yet available to conscious report. Future work combining longitudinal designs and real-time neural measures may further clarify how and when such latent learning signals become accessible for explicit use.